% ================================================================
% 《极小曲面教程》中译本 —— mathtranslation.cls v3.1 之上的兼容层
% 原书: Tobias Holck Colding & William P. Minicozzi II,
%       A Course in Minimal Surfaces, AMS GSM 121, 2011.
% ----------------------------------------------------------------
% 本文件只补齐类文件未提供、而本书确需的部分; 类文件本身不改动.
%   1. 定理类环境与带编号小标题共享同一计数器, 全书按章编号 (1.1, 1.2, ...),
%      与原书版式一致(原书的 "1.1." 既可能是小标题, 也可能是引理).
%   2. \mtpara{标题}     节内带编号的 run-in 小标题(接正文).
%      \mtparasec{标题}  同上, 但独占一行.
%   3. 未编号环境 theorem* / lemma* / ... 与断言 claim.
%   4. 常用记号缩写(曲面、法向、面积、体积等).
% ================================================================

% ---------- 1. 共享计数器 ----------
% 类文件已用 [section] 为每个定理类环境各建一个计数器, 且把它们登记进
% \cl@section 的重置链. 这里只解除环境定义, 保留计数器本身, 以免 \section
% 报 "No counter ... defined".
\makeatletter
\let\theorem\relax\let\endtheorem\relax
\let\definition\relax\let\enddefinition\relax
\let\lemma\relax\let\endlemma\relax
\let\corollary\relax\let\endcorollary\relax
\let\proposition\relax\let\endproposition\relax
\let\remark\relax\let\endremark\relax
\let\example\relax\let\endexample\relax
\let\problem\relax\let\endproblem\relax
\makeatother

% 共享计数器: 按章重置, 显示 "章号.序号".
\newcounter{mtitem}[chapter]
\renewcommand{\themtitem}{\thechapter.\arabic{mtitem}}

\theoremstyle{mtfang}
\newtheorem{theorem}[mtitem]{定理}
\newtheorem{lemma}[mtitem]{引理}
\newtheorem{definition}[mtitem]{定义}
\newtheorem{corollary}[mtitem]{推论}
\newtheorem{proposition}[mtitem]{命题}
\newtheorem{remark}[mtitem]{注}
\newtheorem{conjecture}[mtitem]{猜想}
\theoremstyle{mtsong}
\newtheorem{example}[mtitem]{例}
\newtheorem{problem}[mtitem]{问题}

% 未编号版本(原书中少量不编号的表述).
\newtheorem*{theorem*}{定理}
\newtheorem*{lemma*}{引理}
\newtheorem*{definition*}{定义}
\newtheorem*{corollary*}{推论}
\newtheorem*{proposition*}{命题}
\newtheorem*{remark*}{注}
\newtheorem*{example*}{例}
\newtheorem*{conjecture*}{猜想}

% 断言: 证明内部常用的无编号小块, 自带冒号, 不结束段落.
\newtheorem*{claim*}{断言}
\AtEndEnvironment{example*}{\envenddiamond}
\AtEndEnvironment{remark*}{\envenddiamond}

% ---------- 2. 带编号的小标题 ----------
% 原书的小标题按"节号.序号"编号(第 1 节里的 1.1, 1.2, 1.3; 第 2 节里的
% 2.1, ..., 2.14), 与按章编号的定理类对象(引理 1.1, 定理 1.21, ...)是两套
% 互不相干的计数器 —— 原书里 "1.1" 既可能是小标题也可能是引理, 本译本照此.
\newcounter{mtpara}[section]
\renewcommand{\themtpara}{\thesection.\arabic{mtpara}}

\newcommand{\mtpara}[1]{%
  \par\medskip\noindent\refstepcounter{mtpara}%
  \setcounter{mtparasub}{0}%
  {\headfont\themtpara.\enspace #1.}\quad\ignorespaces
}
\newcommand{\mtparasec}[1]{%
  \par\medskip\noindent\refstepcounter{mtpara}%
  \setcounter{mtparasub}{0}%
  {\headfont\themtpara.\enspace #1.}\par\smallskip\noindent\ignorespaces
}

% 三级小标题: 原书第 8 章 §5.3 下的 5.3.1, 5.3.2, 5.3.3.
\newcounter{mtparasub}
\renewcommand{\themtparasub}{\thesection.\arabic{mtpara}.\arabic{mtparasub}}
\newcommand{\mtparasub}[1]{%
  \par\medskip\noindent\refstepcounter{mtparasub}%
  {\headfont\themtparasub.\enspace #1.}\quad\ignorespaces
}

% ---------- 3. 节标题样式 ----------
% 原书节标题为 "1. The Minimal Surface Equation and Minimal Submanifolds"
% (编号后带句点, 起点左对齐, 无章号前缀).
\renewcommand{\thesection}{\arabic{section}}
\titleformat{\section}[block]
  {\normalfont\large\headfont}
  {\thesection.}
  {0.6em}
  {}
\titlespacing*{\section}{0pt}{2.4ex plus 1ex minus .2ex}{1.0ex plus .2ex}

% ---------- 4. 引用命令 ----------
% 所有定理类环境共用一个计数器, \autoref 无法区分类型, 故按类型分别引用.
\newcommand{\thmref}[1]{定理~\ref{#1}}
\newcommand{\lemref}[1]{引理~\ref{#1}}
\newcommand{\defref}[1]{定义~\ref{#1}}
\newcommand{\propref}[1]{命题~\ref{#1}}
\newcommand{\corref}[1]{推论~\ref{#1}}
\newcommand{\remref}[1]{注~\ref{#1}}
\newcommand{\exref}[1]{例~\ref{#1}}
\newcommand{\secref}[1]{第~\ref{#1}~节}
\newcommand{\chapref}[1]{第~\ref{#1}~章}
\newcommand{\figref}[1]{图~\ref{#1}}

% ---------- 5. 记号缩写 ----------
\newcommand{\R}{\mathbb{R}}
\newcommand{\C}{\mathbb{C}}
\newcommand{\bS}{\mathbf{S}}
\newcommand{\bB}{\mathbf{B}}
\DeclareMathOperator{\Area}{Area}
\DeclareMathOperator{\Vol}{Vol}
\DeclareMathOperator{\Graph}{Graph}
\DeclareMathOperator{\sgn}{sgn}
\DeclareMathOperator{\diam}{diam}
\DeclareMathOperator{\inj}{inj}
\DeclareMathOperator{\osc}{osc}
\DeclareMathOperator{\supp}{supp}
\DeclareMathOperator{\dist}{dist}
\DeclareMathOperator{\Length}{Length}
\DeclareMathOperator{\Hess}{Hess}
\newcommand{\tub}{\mathcal{T}}
\DeclareMathOperator{\Span}{span}
\newcommand{\dive}{\operatorname{div}}
\newcommand{\ii}{\mathrm{i}}

% ---------- 6. 图表标题 ----------
\usepackage{caption}
\captionsetup{font=small,labelsep=period,skip=6pt}

% 原书的编号列表一律用 (1)(2)... / (i)(ii)... 的形式, 这里统一.
\usepackage{enumitem}
\setlist[enumerate]{label=(\arabic*),leftmargin=2.4em,itemsep=2pt,topsep=2pt}

% 插图统一入口: \mtfigure[宽度比例]{文件名}{图题}{标签}
\newcommand{\mtfigure}[4][0.45]{%
  \begin{figure}[htbp]
    \centering
    \includegraphics[width=#1\textwidth]{fig/#2}%
    \caption{#3}%
    \label{#4}%
  \end{figure}}

% 并排两图(原书 1.27/1.28、1.31/1.32、1.33/1.34 即为左右并置)
\newcommand{\mtfigurepair}[6]{%
  \begin{figure}[htbp]
    \centering
    \begin{minipage}[t]{0.485\textwidth}
      \centering
      \includegraphics[width=\linewidth,height=0.16\textheight,keepaspectratio]{fig/#1}%
      \caption{#2}\label{#3}
    \end{minipage}\hfill
    \begin{minipage}[t]{0.485\textwidth}
      \centering
      \includegraphics[width=\linewidth,height=0.16\textheight,keepaspectratio]{fig/#4}%
      \caption{#5}\label{#6}
    \end{minipage}
  \end{figure}}

% ---------- 7. 术语索引改为不编号章节(双栏) ----------
% 类文件的 \printterminology 假定文档处于 \appendix 状态, 用 \chapter 排
% "术语索引"(于是显示为"附录 X 术语索引"). 本书没有附录, 术语索引直接跟在
% 参考文献之后, 故这里自建一份: 章节用 \chapter*(带目录条目).
%
% 词条数超过 50 后单栏一页放不下, 改为双栏索引表. 关键在于 \MTTermList
% 本身就是一串 \MTTermRow{中文}{英文}{标签} 调用, 所以**不必遍历、也不必
% 拆分**: 把 \MTTermRow 重定义为一段"术语 …… 页码", 再直接执行 \MTTermList
% 即可按登记顺序逐条输出; multicols 负责分栏与均衡, 阅读顺序天然是
% "先自上而下、再从左到右"。
%
% (曾试过先切成 expl3 序列再按下标两两并排, 但 \MTTermRow 不可展开,
%  \MTTermRow \seq_item:Nn ... 会把 \seq_item:Nn 当成第一个参数吞掉;
%  即使加 \exp_after:wN 也不稳, 故弃用。)
\usepackage{multicol}
\makeatletter
\newcommand{\printterminologyunnumbered}{%
  \cleardoublepage
  \phantomsection
  \chapter*{术语索引}%
  \addcontentsline{toc}{chapter}{术语索引}%
  \markboth{术语索引}{术语索引}%
  \begingroup
  \scriptsize
  \setlength{\parindent}{0pt}%
  \setlength{\parskip}{1.1pt}%
  \renewcommand{\MTTermRow}[3]{%
    \par\noindent\hangindent=1.1em\hangafter=1
    {\termfont ##1}~##2\nobreak\dotfill\nobreak
    \hyperref[##3]{\pageref*{##3}}\par}%
  \begin{multicols}{2}
    \raggedright
    \MTTermList
  \end{multicols}
  \endgroup}
\makeatother
